\chapter{Im lặng và suy nghĩ}
\markboth{Im lặng và suy nghĩ}{Silence and Thinking}


\section{Không phải ai nói nhiều cũng là người hiểu nhiều.}
\begin{truyen}{Không phải ai nói nhiều cũng là người hiểu nhiều.}{Không phải ai nói nhiều cũng là người hiểu nhiều.}
\chuhoa{C}{âu thầy nói câu ấy. Cả lớp im lặng — không ai nói, mỗi người đang nghĩ về điều vừa nghe.}
\textit{(The teacher said that. The class fell silent — no one spoke; everyone was thinking about what they had just heard.)}
\end{truyen}

\begin{baihoc}
Im lặng sau câu nói đúng lúc là dấu hiệu của tư duy.
\textit{(Silence after the right sentence is a sign of thinking.)}
\end{baihoc}

\begin{ghinhoanh}
\textbf{Short English to remember:}

Some sentences make the whole class think in silence.

\end{ghinhoanh}

\begin{tuhocnhanh}
1. Câu này khiến em nghĩ đến điều gì? \textit{What does this make you think of?}

2. Em có thể dùng câu này trong tình huống nào? \textit{When could you use this sentence?}
\end{tuhocnhanh}
\ngancach


\section{Im lặng đôi khi là câu trả lời hay nhất.}
\begin{truyen}{Im lặng đôi khi là câu trả lời hay nhất.}{Im lặng đôi khi là câu trả lời hay nhất.}
\chuhoa{C}{âu thầy nói câu ấy. Cả lớp im lặng — không ai nói, mỗi người đang nghĩ về điều vừa nghe.}
\textit{(The teacher said that. The class fell silent — no one spoke; everyone was thinking about what they had just heard.)}
\end{truyen}

\begin{baihoc}
Im lặng sau câu nói đúng lúc là dấu hiệu của tư duy.
\textit{(Silence after the right sentence is a sign of thinking.)}
\end{baihoc}

\begin{ghinhoanh}
\textbf{Short English to remember:}

Some sentences make the whole class think in silence.

\end{ghinhoanh}

\begin{tuhocnhanh}
1. Câu này khiến em nghĩ đến điều gì? \textit{What does this make you think of?}

2. Em có thể dùng câu này trong tình huống nào? \textit{When could you use this sentence?}
\end{tuhocnhanh}
\ngancach


\section{Em đang chạy theo điểm hay đang chạy theo hiểu?}
\begin{truyen}{Em đang chạy theo điểm hay đang chạy theo hiểu?}{Em đang chạy theo điểm hay đang chạy theo hiểu?}
\chuhoa{C}{âu thầy nói câu ấy. Cả lớp im lặng — không ai nói, mỗi người đang nghĩ về điều vừa nghe.}
\textit{(The teacher said that. The class fell silent — no one spoke; everyone was thinking about what they had just heard.)}
\end{truyen}

\begin{baihoc}
Im lặng sau câu nói đúng lúc là dấu hiệu của tư duy.
\textit{(Silence after the right sentence is a sign of thinking.)}
\end{baihoc}

\begin{ghinhoanh}
\textbf{Short English to remember:}

Some sentences make the whole class think in silence.

\end{ghinhoanh}

\begin{tuhocnhanh}
1. Câu này khiến em nghĩ đến điều gì? \textit{What does this make you think of?}

2. Em có thể dùng câu này trong tình huống nào? \textit{When could you use this sentence?}
\end{tuhocnhanh}
\ngancach


\section{Câu hỏi quan trọng hơn câu trả lời — vì nó mở đường.}
\begin{truyen}{Câu hỏi quan trọng hơn câu trả lời — vì nó mở đường.}{Câu hỏi quan trọng hơn câu trả lời — vì nó mở đườn}
\chuhoa{C}{âu thầy nói câu ấy. Cả lớp im lặng — không ai nói, mỗi người đang nghĩ về điều vừa nghe.}
\textit{(The teacher said that. The class fell silent — no one spoke; everyone was thinking about what they had just heard.)}
\end{truyen}

\begin{baihoc}
Im lặng sau câu nói đúng lúc là dấu hiệu của tư duy.
\textit{(Silence after the right sentence is a sign of thinking.)}
\end{baihoc}

\begin{ghinhoanh}
\textbf{Short English to remember:}

Some sentences make the whole class think in silence.

\end{ghinhoanh}

\begin{tuhocnhanh}
1. Câu này khiến em nghĩ đến điều gì? \textit{What does this make you think of?}

2. Em có thể dùng câu này trong tình huống nào? \textit{When could you use this sentence?}
\end{tuhocnhanh}
\ngancach


\section{Người thông minh nhất trong phòng thường là người im lặng lắng nghe.}
\begin{truyen}{Người thông minh nhất trong phòng thường là người im lặng lắ...}{Người thông minh nhất trong phòng thường là người }
\chuhoa{C}{âu thầy nói câu ấy. Cả lớp im lặng — không ai nói, mỗi người đang nghĩ về điều vừa nghe.}
\textit{(The teacher said that. The class fell silent — no one spoke; everyone was thinking about what they had just heard.)}
\end{truyen}

\begin{baihoc}
Im lặng sau câu nói đúng lúc là dấu hiệu của tư duy.
\textit{(Silence after the right sentence is a sign of thinking.)}
\end{baihoc}

\begin{ghinhoanh}
\textbf{Short English to remember:}

Some sentences make the whole class think in silence.

\end{ghinhoanh}

\begin{tuhocnhanh}
1. Câu này khiến em nghĩ đến điều gì? \textit{What does this make you think of?}

2. Em có thể dùng câu này trong tình huống nào? \textit{When could you use this sentence?}
\end{tuhocnhanh}
\ngancach


\section{Sách không thay em nghĩ; sách mời em nghĩ.}
\begin{truyen}{Sách không thay em nghĩ; sách mời em nghĩ.}{Sách không thay em nghĩ; sách mời em nghĩ.}
\chuhoa{C}{âu thầy nói câu ấy. Cả lớp im lặng — không ai nói, mỗi người đang nghĩ về điều vừa nghe.}
\textit{(The teacher said that. The class fell silent — no one spoke; everyone was thinking about what they had just heard.)}
\end{truyen}

\begin{baihoc}
Im lặng sau câu nói đúng lúc là dấu hiệu của tư duy.
\textit{(Silence after the right sentence is a sign of thinking.)}
\end{baihoc}

\begin{ghinhoanh}
\textbf{Short English to remember:}

Some sentences make the whole class think in silence.

\end{ghinhoanh}

\begin{tuhocnhanh}
1. Câu này khiến em nghĩ đến điều gì? \textit{What does this make you think of?}

2. Em có thể dùng câu này trong tình huống nào? \textit{When could you use this sentence?}
\end{tuhocnhanh}
\ngancach


\section{Nếu em không đặt câu hỏi, em đang tin tất cả.}
\begin{truyen}{Nếu em không đặt câu hỏi, em đang tin tất cả.}{Nếu em không đặt câu hỏi, em đang tin tất cả.}
\chuhoa{C}{âu thầy nói câu ấy. Cả lớp im lặng — không ai nói, mỗi người đang nghĩ về điều vừa nghe.}
\textit{(The teacher said that. The class fell silent — no one spoke; everyone was thinking about what they had just heard.)}
\end{truyen}

\begin{baihoc}
Im lặng sau câu nói đúng lúc là dấu hiệu của tư duy.
\textit{(Silence after the right sentence is a sign of thinking.)}
\end{baihoc}

\begin{ghinhoanh}
\textbf{Short English to remember:}

Some sentences make the whole class think in silence.

\end{ghinhoanh}

\begin{tuhocnhanh}
1. Câu này khiến em nghĩ đến điều gì? \textit{What does this make you think of?}

2. Em có thể dùng câu này trong tình huống nào? \textit{When could you use this sentence?}
\end{tuhocnhanh}
\ngancach


\section{Điều em sợ nói ra thường là điều em cần nói nhất.}
\begin{truyen}{Điều em sợ nói ra thường là điều em cần nói nhất.}{Điều em sợ nói ra thường là điều em cần nói nhất.}
\chuhoa{C}{âu thầy nói câu ấy. Cả lớp im lặng — không ai nói, mỗi người đang nghĩ về điều vừa nghe.}
\textit{(The teacher said that. The class fell silent — no one spoke; everyone was thinking about what they had just heard.)}
\end{truyen}

\begin{baihoc}
Im lặng sau câu nói đúng lúc là dấu hiệu của tư duy.
\textit{(Silence after the right sentence is a sign of thinking.)}
\end{baihoc}

\begin{ghinhoanh}
\textbf{Short English to remember:}

Some sentences make the whole class think in silence.

\end{ghinhoanh}

\begin{tuhocnhanh}
1. Câu này khiến em nghĩ đến điều gì? \textit{What does this make you think of?}

2. Em có thể dùng câu này trong tình huống nào? \textit{When could you use this sentence?}
\end{tuhocnhanh}
\ngancach


\section{Học không phải để trả lời đúng — học để biết hỏi đúng.}
\begin{truyen}{Học không phải để trả lời đúng — học để biết hỏi đúng.}{Học không phải để trả lời đúng — học để biết hỏi đ}
\chuhoa{C}{âu thầy nói câu ấy. Cả lớp im lặng — không ai nói, mỗi người đang nghĩ về điều vừa nghe.}
\textit{(The teacher said that. The class fell silent — no one spoke; everyone was thinking about what they had just heard.)}
\end{truyen}

\begin{baihoc}
Im lặng sau câu nói đúng lúc là dấu hiệu của tư duy.
\textit{(Silence after the right sentence is a sign of thinking.)}
\end{baihoc}

\begin{ghinhoanh}
\textbf{Short English to remember:}

Some sentences make the whole class think in silence.

\end{ghinhoanh}

\begin{tuhocnhanh}
1. Câu này khiến em nghĩ đến điều gì? \textit{What does this make you think of?}

2. Em có thể dùng câu này trong tình huống nào? \textit{When could you use this sentence?}
\end{tuhocnhanh}
\ngancach


\section{Một phút im lặng suy nghĩ đáng giá hơn mười phút nói mà không nghĩ.}
\begin{truyen}{Một phút im lặng suy nghĩ đáng giá hơn mười phút nói mà khôn...}{Một phút im lặng suy nghĩ đáng giá hơn mười phút n}
\chuhoa{C}{âu thầy nói câu ấy. Cả lớp im lặng — không ai nói, mỗi người đang nghĩ về điều vừa nghe.}
\textit{(The teacher said that. The class fell silent — no one spoke; everyone was thinking about what they had just heard.)}
\end{truyen}

\begin{baihoc}
Im lặng sau câu nói đúng lúc là dấu hiệu của tư duy.
\textit{(Silence after the right sentence is a sign of thinking.)}
\end{baihoc}

\begin{ghinhoanh}
\textbf{Short English to remember:}

Some sentences make the whole class think in silence.

\end{ghinhoanh}

\begin{tuhocnhanh}
1. Câu này khiến em nghĩ đến điều gì? \textit{What does this make you think of?}

2. Em có thể dùng câu này trong tình huống nào? \textit{When could you use this sentence?}
\end{tuhocnhanh}
\ngancach
